\documentclass{controle}
\usepackage{main}

\title{Évaluation n°11 : Introductions aux vecteurs du plan}
\date{30 Janvier 2026}
\author{Seconde 3}

% \printanswers

\begin{document}
\titleandname{}
\version

\begin{questions}
\question[1,5] Citer les trois caractéristiques d'un vecteur.

\begin{solution}[0.3cm]
Les trois caractéristiques d'un vecteur sont :
\begin{itemize}
\item sa direction;
\item son sens;
\item sa norme
\end{itemize}
\end{solution}

\question[1] Rappeler la relation de Chasles entre les vecteurs $\vect{PQ}$ et $\vect{QR}$.

\begin{solution}[0.3cm]
\begin{equation*}
\vect{PQ} + \vect{QR} = \vect{AC} 
\end{equation*}
\end{solution}

\question[7,5] Répondre aux questions suivantes. Tous les triangles représentés sur la figure sont équilatérals.

\begin{center}
\begin{tikzpicture}
\coordinate (A) at (0,0);
\coordinate (B) at (0:2);
\coordinate (C) at (60:2);
\coordinate (D) at ($(B) + (60:2)$);
\coordinate (E) at (0:4);
\coordinate (F) at (60:4);

\draw (A) node[below left] {$A$} -- (B) node[below] {$B$} 
  -- (C) node[left] {$C$} -- cycle;
\draw (B) -- (D) node[right] {$D$} -- (E) node[below right] {$E$} -- cycle;
\draw (C) -- (D) -- (F) node[above] {$F$} -- cycle;
\end{tikzpicture}
\end{center}
\begin{parts}
\part Donner un vecteur égal à $\vect{AB}$.
\part Donner le représentant de $\vect{DB}$ ayant pour origine $F$.
\part Donner le représentant de $\vect{CF}$ ayant pour extrémité $C$.
\part Donner un vecteur de même direction que $\vect{EB}$, mais de sens opposé.
\part Donner un vecteur de même direction et de même sens que $\vect{DB}$ mais de norme différente.
\end{parts}

\begin{solution}
\begin{parts}
\part $\vect{CD}$
\part $\vect{FC}$
\part $\vect{AC}$
\part $\vect{BE}$
\part $\vect{FA}$
\end{parts}
\end{solution}

\end{questions}

\newpage
\setcounter{page}{1}
\titleandname{}
\version

\begin{questions}
\question[1,5] Citer les trois caractéristiques d'un vecteur.

\begin{solution}[0.3cm]
Les trois caractéristiques d'un vecteur sont :
\begin{itemize}
\item sa direction;
\item son sens;
\item sa norme
\end{itemize}
\end{solution}

\question[1] Rappeler la relation de Chasles entre les vecteurs $\vect{MN}$ et $\vect{NO}$.

\begin{solution}[0.3cm]
\begin{equation*}
\vect{MN} + \vect{NO} = \vect{MO} 
\end{equation*}
\end{solution}

\question[7,5] Répondre aux questions suivantes. Tous les triangles représentés sur la figure sont équilatérals.

\begin{center}
\begin{tikzpicture}
\coordinate (A) at (0,0);
\coordinate (B) at (0:2);
\coordinate (C) at (-60:2);
\coordinate (D) at ($(B) + (-60:2)$);
\coordinate (E) at (0:4);
\coordinate (F) at (-60:4);

\draw (A) node[above left] {$A$} -- (B) node[above] {$B$} 
  -- (C) node[left] {$C$} -- cycle;
\draw (B) -- (D) node[right] {$D$} -- (E) node[above right] {$E$} -- cycle;
\draw (C) -- (D) -- (F) node[below] {$F$} -- cycle;
\end{tikzpicture}
\end{center}
\begin{parts}
\part Donner un vecteur égal à $\vect{BC}$.
\part Donner le représentant de $\vect{DB}$ ayant pour extrémité $C$.
\part Donner le représentant de $\vect{AB}$ ayant pour origine $B$.
\part Donner un vecteur de même direction que $\vect{FD}$, mais de sens opposé.
\part Donner un vecteur de même direction et de même sens que $\vect{DC}$ mais de norme différente.
\end{parts}

\begin{solution}
\begin{parts}
\part $\vect{ED}$
\part $\vect{FC}$
\part $\vect{BE}$
\part $\vect{BC}$
\part $\vect{EA}$
\end{parts}
\end{solution}


\end{questions}
\end{document}